\originalchapter{干涉、衍射与分辨率}\label{ch:diffraction}
\originalsection{相干叠加与薄膜颜色}
两个波在同一区域存在时, 总电场是场的和, 而探测器常记录平方的时间平均. 同频、同偏振、相位关系稳定的两束光给
\[
I=I_1+I_2+2\sqrt{I_1I_2}\cos\delta.
\]
$\delta$ 是两振幅在探测点的相位差. 对一般单位 Jones 向量 $\boldsymbol u_1,\boldsymbol u_2$, 若 $\delta$ 定义为第二束相对第一束的整体相位, 则交叉项为 $2\sqrt{I_1I_2}\Rea[\boldsymbol u_1^\dagger\boldsymbol u_2\e^{\ii\delta}]$. 实线偏振退化为实点积乘 $\cos\delta$, 正交偏振使交叉项为零; 若相位随时间随机变化, 长时间平均交叉项可消失. 相干性是能看到稳定条纹的条件, 不能对任何两束光都画出固定明暗纹. 局部暗纹不表示能量消失, 相邻方向或另一输出端的能量增强, 总能量需按整个系统核对.

肥皂膜的颜色来自前、后表面反射波的相位差. 设膜折射率 $n_f$, 两侧空气相同, 厚度 $d$, 正入射真空波长 $\lambda_0$. 前面由低折射率到高折射率反射有 $\pi$ 相移, 背面由高到低反射无相移. 两束回到同一侧的相对相位可写为 $\delta=4\pi n_fd/\lambda_0+\pi$ 模 $2\pi$. 反射增强为 $2n_fd=(m+1/2)\lambda_0$, 反射减弱为 $2n_fd=m\lambda_0$; $m$ 取使厚度非负的整数. 斜入射把传播项改为 $4\pi n_fd\cos\theta_f/\lambda_0$, 膜内角由折射定律求出.

两束近似足以说明颜色, 但严格零反射还涉及两束幅度及多次内部反射. 正入射的完整无损薄层振幅为
\[
r_{\rm film}=\frac{r_{01}+r_{12}\e^{2\ii k_fd}}
{1+r_{01}r_{12}\e^{2\ii k_fd}},\qquad
r_{ab}=\frac{n_a-n_b}{n_a+n_b},\quad k_f=2\pi n_f/\lambda_0.
\]
令 $t_{ab}=2n_a/(n_a+n_b)$ 为单界面电场透射系数. 首次进入膜后反射返回的振幅为 $t_{01}t_{10}r_{12}\e^{2\ii k_fd}$, 每增加一次往返乘 $r_{10}r_{12}\e^{2\ii k_fd}$. 因此总反射为
\[r_{01}+t_{01}t_{10}r_{12}\e^{2\ii k_fd}\sum_{j=0}^{\infty}(r_{10}r_{12}\e^{2\ii k_fd})^j.\]
有限正折射率无损界面有 $|r_{10}r_{12}|<1$, 几何级数收敛. 使用 $r_{10}=-r_{01}$、$t_{01}t_{10}=1-r_{01}^2$, 整理便得上式, 分母包含多次反射. 两侧相同且 $d\to0$ 时 $r_{12}=-r_{01}$, 得零反射, 比只检查相位更完整. 若膜在基底上, 两反射的相移要重新判断; 不能直接照抄肥皂膜明暗条件. $n_f=1.5,d=100\,\mathrm{nm}$ 时首个反射增强约在 $600\,\mathrm{nm}$; 更厚的膜在可见波段有许多密集级次, 有限光源带宽与空间分辨率会使颜色平均.

\originalsection{双缝与相控阵}
两等幅相干小孔间距 $d$, 在远处角度 $\theta$ 的光程差近似 $d\sin\theta$. 条件是观察距离远大于孔间距, 更严格的夫琅禾费条件将在下节说明. 相位差 $\delta=kd\sin\theta$, $k=2\pi/\lambda$, 则 $I=4I_s\cos^2(\delta/2)$, $I_s$ 是单孔在该方向的强度. 亮纹 $d\sin\theta=m\lambda$, 暗纹 $d\sin\theta=(m+1/2)\lambda$. 小角屏距 $L$ 时条纹间距约为 $\lambda L/d$.

$N$ 个等间距源, 相邻源本身增加相位 $\phi$, 总相位间隔为 $\delta=kd\sin\theta-\phi$. 几何级数求和给
\[
\sum_{j=0}^{N-1}\e^{-\ii j\delta}
=\e^{-\ii(N-1)\delta/2}\frac{\sin(N\delta/2)}{\sin(\delta/2)},\qquad
I=I_s\left[\frac{\sin(N\delta/2)}{\sin(\delta/2)}\right]^2.
\]
分母为零时需取极限, 主极大强度为 $N^2I_s$. 相邻主极大之间, $\delta=2\pi l/N$, $l=1,\ldots,N-1$ 给零点, 内部有 $N-2$ 个次极大. 主极大邻近的第一零点相位偏差 $2\pi/N$, 角宽约 $\lambda/(Nd\cos\theta_0)$. 调整 $\phi$ 可令主束满足 $kd\sin\theta_0=\phi+2\pi m$, 是相控雷达转向的原理. 大间距也可能出现多个满足条件的主束, 不能只说增加 $d$ 永远改善唯一方向的分辨能力.

\originalsection{单缝的连续叠加}
孔径上的各位置携带不同传播相位, 衍射就是对连续分布振幅求和. 对均匀照明、宽度 $a$ 的单缝, 远场复振幅正比于
\[
\int_{-a/2}^{a/2}\e^{-\ii kx\sin\theta}\dd x
=a\frac{\sin\beta}{\beta},\qquad \beta=\frac{\pi a\sin\theta}{\lambda}.
\]
故 $I=I(0)\sinc^2\beta$, $\sinc\beta=\sin\beta/\beta$, $\sinc0=1$. 除中心以外的零点满足 $a\sin\theta=m\lambda$, $m$ 为非零整数. 中央亮纹比邻近亮纹宽, 强度包络按 $1/\beta^2$ 衰减, 次极大的精确位置由 $\tan\beta=\beta$ 求出, 不能说它们全部等间隔.

远场近似把从孔径点到探测点的距离只保留线性相位, 要求被忽略的二次相位 $ka^2/L$ 足够小, 常写 $L\gg a^2/\lambda$. 也可用透镜在焦平面实现相应的傅里叶分布. 近场不满足这个条件时, 要保留二次相位, 属于菲涅耳衍射. 双缝有限缝宽 $a$, 中心距 $d$ 的强度为 $4I_s(0)\sinc^2\beta\cos^2(\delta/2)$; $N$ 缝则是单缝包络乘上一节的阵列因子. 若干主极大落在单缝零点上, 会出现缺级.

\begin{center}
\begin{tikzpicture}
\begin{axis}[width=.76\linewidth,height=4.5cm,xlabel={$\beta$},ylabel={$I/I(0)$},xmin=-10,xmax=10,ymin=0,ymax=1.08,samples=201,domain=-10:10]
\addplot[thick,structurecolor]{abs(x)<.0001 ? 1 : (sin(deg(x))/x)^2};
\end{axis}
\end{tikzpicture}\par
单缝夫琅禾费衍射; $\beta$ 用弧度表示, 绘图时显式换成三角函数的角度单位.
\end{center}

\originalsection{二维孔径与圆孔分辨率}
二维孔径透过函数 $A(x,y)$ 的远场振幅为其二维傅里叶积分 $\int\int A(x,y)\e^{-\ii(k_xx+k_yy)}\dd x\dd y$, 其中 $k_x,k_y$ 由观察方向决定. 矩形孔 $a\times b$ 给两个 $\sinc$ 的乘积. 圆孔半径 $R=D/2$ 的积分用极坐标完成, 记横向波数 $q=k\sin\theta$:
\[
\int_0^Rr\dd r\int_0^{2\pi}\e^{-\ii qr\cos\varphi}\dd\varphi
=2\pi\int_0^RrJ_0(qr)\dd r
=\frac{2\pi R}{q}J_1(qR).
\]
这里 $J_0(z)$ 可由角积分定义, $J_1$ 是第一类1阶贝塞尔函数; $[zJ_1(z)]'=zJ_0(z)$ 给最后一步. 从幂级数 $J_1(z)=z/2+O(z^3)$ 可知中心振幅为 $\pi R^2$. 因而正确归一强度为
\[
\frac{I(\theta)}{I(0)}=\left[\frac{2J_1(\beta)}{\beta}\right]^2,
\quad \beta=\frac{\pi D\sin\theta}{\lambda}.
\]
$J_1$ 的第一个正零点约 $3.8317$, 第一暗环满足 $\sin\theta\simeq1.22\lambda/D$. 数值零点是特定特殊函数的计算结果, 不是所有衍射问题的通用证明.

两个互不相干点源分别产生艾里斑. 瑞利判据约定一个斑中心位于另一个斑的第一暗环时可辨, 所以小角分辨率约 $1.22\lambda/D$; 这是基于特定点扩散函数的分辨约定, 还受噪声与探测器影响. 对 $D=5\,\mathrm{mm},\lambda=500\,\mathrm{nm}$, 得 $1.22\times10^{-4}\,\mathrm{rad}$. 对 $D=1\,\mathrm{mm}$ 的均匀圆孔, 在 $4\times10^8\,\mathrm m$ 处第一暗环半径约 $244\,\mathrm{km}$, 说明有限孔径不能生成永久不扩张的光束. 此圆孔系数不能不加说明用于单缝, 也不等同于高斯激光束的发散定义.

衍射光栅在第 $m$ 级主极大满足 $d\sin\theta=m\lambda$. 微分得两近波长角距 $\Delta\theta\simeq m\Delta\lambda/(d\cos\theta)$. 与阵列主峰第一零点角宽比较, 瑞利分辨本领为 $\lambda/\Delta\lambda\simeq mN$. 此式取正级次 $m\ge1$（负级次以 $|m|$ 计）; 零级不能用此式分辨波长. 它假定照亮 $N$ 条等效缝且包络没有抑制所用级次, 给出“更多相干源导致更窄主峰”的可量化光谱意义.
